% Generated draft: requires mathematical and source review.
\documentclass[11pt]{article}
\usepackage[margin=25mm]{geometry}
\usepackage{fontspec}
\setmainfont{Latin Modern Roman}
\usepackage{amsmath,amssymb,mathtools,unicode-math}
\setmathfont{Latin Modern Math}
\usepackage{enumitem,needspace,xurl,hyperref}
\setlength{\parindent}{0pt}
\setlength{\parskip}{5pt}
\setlength{\emergencystretch}{3em}
\newcommand{\N}{\mathbb{N}}
\newcommand{\Z}{\mathbb{Z}}
\newcommand{\Q}{\mathbb{Q}}
\newcommand{\R}{\mathbb{R}}
\newcommand{\C}{\mathbb{C}}
\newcommand{\F}{\mathbb{F}}
\newcommand{\A}{\mathbb{A}}
\newcommand{\PP}{\mathbb{P}}
\newcommand{\E}{\mathbb{E}}
\newcommand{\eps}{\varepsilon}
\newcommand{\dd}{\,\mathrm{d}}
\begin{document}
\section*{UnsolvedMath ID: 30004357}
\textbf{Source problem code:} catalogue entry
\subsection{Definitions and mathematical statement}
Let \(k\geq 3\) and \(n\geq 1\) be integers, and write \([k]=\{1,\ldots,k\}\). A family \(A\subseteq[k]^n\) is \emph{intersecting} if
\[
\forall x,y\in A\quad \exists i\in\{1,\ldots,n\}\quad x_i=y_i.
\]
For a precise coordinate-permutation interpretation of symmetry, let \(S_n\) act on \([k]^n\) by \((\sigma\cdot x)_i=x_{\sigma^{-1}(i)}\), and define
\[
G_A=\{\sigma\in S_n:\sigma\cdot A=A\}.
\]
Say that \(A\) is coordinate-transitive if
\[
\forall i,j\in\{1,\ldots,n\}\quad \exists\sigma\in G_A\quad \sigma(i)=j.
\]
Under this interpretation, define
\[
M_k(n)=\max\bigl\{|A|:A\subseteq[k]^n\text{ is intersecting and coordinate-transitive}\bigr\}.
\]
The explicit extremal question is to determine sharp upper bounds for \(M_k(n)\), including its asymptotic growth as \(n\to\infty\) for each fixed integer \(k\geq3\).

The broader source question asks for bounds after excluding highly asymmetric, single-coordinate extremal examples. A single-coordinate family has the form
\[
D_{i,a}=\{x\in[k]^n:x_i=a\},\qquad i\in\{1,\ldots,n\},\quad a\in[k].
\]
The source does not specify whether exclusion means forbidding equality with such a family, forbidding containment in one, or imposing a quantitative condition on resemblance to one. Thus this broader part does not yet define a unique extremal problem.

The definition of \(G_A\) above specifies one interpretation of the transitive subquestion. The source does not state whether its automorphisms are restricted to coordinate permutations or may also permute alphabet symbols within coordinates; the latter convention could yield a different admissible class. The source also leaves the range of \(n\) implicit. The convention \(n\geq1\) makes the definitions explicit, but at \(n=1\) coordinate transitivity is automatic and does not exclude single-coordinate families. These qualifications require mathematical review.
\subsection{Short English statement}
An intersecting family consists of words such that every pair agrees in some coordinate. The explicit subquestion asks how large such a family can be when its coordinate symmetries act transitively. The intended automorphism convention and the broader rule excluding single-coordinate examples remain unspecified. The supplied dated curation labels the problem partially solved.
\subsection{Sources}
\begin{itemize}
\item Supplied local UnsolvedMath registry record, permanent ID \(30004357\), title \emph{Symmetric Intersecting Families in Discrete Product Spaces}, exact source code \texttt{catalogue entry}, in the Oberwolfach Reports Open Problems collection. The supplied proposed year is \(2020\); no individual proposer is identified.
\item Supplied report citation: \emph{Combinatorics (2021)}, Oberwolfach Reports, DOI \texttt{10.4171/owr/2020/1}, pp.\ 6--89. \url{https://doi.org/10.4171/owr/2020/1}. The supplied background describes these reports as documenting mathematical workshops and open-problem discussions.
\item Supplied literature reference: Eberhard--Kahn--Narayanan--Spirkl, \url{https://arxiv.org/abs/1909.11578}. No title, theorem number, or precise correspondence between its symmetry convention and this record is supplied.
\item Literature assessment dated \texttt{2026-08-22}: the supplied curation states that, for fixed alphabet size at least three, symmetric intersecting families have size \(o(k^n)\), while the broader extremal classification after excluding one-coordinate examples remains incomplete. The supplied status is \texttt{partially\_solved}. This is supplied dated curation derived from the OWR material, not an independently verified literature assessment; its time-sensitive status and cited claims require independent verification.
\item The registry labels the statement \texttt{corrected\_verified} and says its mathematical content was preserved while notation, typography, omitted definitions, or duplicate clauses were repaired. This is supplied editorial metadata, not independent mathematical verification of this draft.
\end{itemize}
\end{document}
